\documentclass{article}
\usepackage{graphicx} % Required for inserting images

\title{Sutton and Barton Exercises Chapter 4}
\author{danieltocco0 }
\date{October 2025}
\usepackage{imakeidx}
\makeindex
\usepackage{amsmath}
\begin{document}
\section{Chapter 4}
\textit{Exercise 4.1} In Example 4.1, if $\pi$ is the equiprobable random policy, what is $q_{\pi}(11 , down )$?
What is $q_{\pi}(7, down)?$ \\
\\

For $q_{\pi}(11 , down )$, we transition to S' = $S^{+}$ with reward -1 and a future 

expected reward of 0. Thus, q $q_{\pi}(11 , down )$ = -1.\\

For $q_{\pi}(7, down)$, we transition to S' = 11 with reward -1 and a future

expected reward of $-\gamma14$. $\gamma$ is just 1. Thus, $q_{\pi}(7, down)? = -15.$ \\
\\\\
\textit{Exercise 4.2} Suppose a new state 15 is added to the gridworld just below
state 13, and its actions, left, up, right, and down, take the agent to states 12, 13, 14,
and 15, respectively. Assume that the transitions from the original states are unchanged.
What, then, is $v_{\pi}(15)$ for the equiprobable random policy? Now suppose the dynamics of
state 13 are also changed, such that action down from state 13 takes the agent to the new
state 15. What is $v_{\pi}(15)$ for the equiprobable random policy in this case?\\
\\

$v_{\pi}(15)$ = -1 + $\frac{1}{4}(-22 - 20 - 14 + v_{\pi}(15))$ $=>$

$v_{\pi}(15)$ = $\frac{4}{3}(-1 - 14) = -20$. \\

For part two of this question, we need to look at $v_{\pi}(13)$ and $v_{\pi}(15)$: 

$v_{\pi}(13)$ = -1 + $\frac{1}{4}(-22 - 20 - 14 + v_{\pi}(15)) =>$

$v_{\pi}(13) - \frac{1}{4}v_{\pi}(15) = -15\\

and from the first part of the question we have:

$\frac{3}{4}v_{\pi}(15) = -1 + \frac{1}{4}(-22 + v_{\pi}(13) - 14) =>$

$\frac{3}{4}v_{\pi}(15) - \frac{1}{4}v_{\pi}(13) = -10.\\

This is a system of linear equations, and doing the math out:\\

$\frac{4}{5}v_{\pi}(15) - 4 - \frac{1}{4}v_{\pi}(15) = -15 =>$

$v_{\pi}(15) = -20$.

Intuitively, this makes sense as the dynamics and returns are the same from 

each state $v_{\pi}(9)$ also yields -20; you can even simply retrace the arrows on 

the right hand side under "optimal policy" and the arrows in terms of 

direction and steps until termination are the same.\\
\\\\
\textit{Exercise 4.3} What are the equations analogous to (4.3), (4.4), and (4.5) for the action-value function $q_{\pi}$ and its successive approximation by a sequence of functions $q_{0}, q_{1}, q_{2}, . . . ?
\\

4.3) $q_{\pi}(s, a) = \mathbb{E}[R_{t+1} + \Sigma_{a_{t+1}}\pi(a_{t+1}|s_{t+1})\gamma q_{\pi}(s_{t+1}, a_{t+1})|s, a]$

4.4) $q_{\pi}(s, a) = \Sigma_{s', r'}p(s', r|a, s)[r + \Sigma_{a'}\pi(a'|s')\gamma q_{\pi}(s', a')]$

4.5) $q_{k+1}(s, a) = \Sigma_{s', r'}p(s', r|a, s)[r + \Sigma_{a'}\pi(a'|s')\gamma q_{k}(s', a')]$
\\\\
\textit{Exercise 4.4} The policy iteration algorithm on page 80 has a subtle bug in that it may
never terminate if the policy continually switches between two or more policies that are
equally good. This is ok for pedagogy, but not for actual use. Modify the pseudocode so
that convergence is guaranteed. 
\\

I think there are a few ways to fix this. One way would to see if the policy 

values themselves are within some some small threshold $\theta_{thresh}$ 

despite following different actions. If they are within the threshold, 

then just set policy-stable = true. The other way I am thinking is to 

look back two actions. That is, since we are seeing if action a is better 

than a-1, then we can see if action a-2 is equal to action a, and if yes, 

we know we are looping between these two actions and can set 

policy-stable = true.
\\\\
\textit{Exercise 4.5} How would policy iteration be defined for action values? Give a complete
algorithm for computing $q_{*}$, analogous to that on page 80 for computing $v_{*}$. Please pay
special attention to this exercise, because the ideas involved will be used throughout the
rest of the book. 
\\

1. Policy is arbitrarily initiated as before\\

2. Policy Evaluation\\

\indent Loop:\\
\indent \indent $\Delta \leftarrow 0$\\
\indent \indent Loop for each s $\in$ S:\\
\indent \indent \indent $\textit{q} \leftarrow Q(s,a)$\\
\indent \indent \indent $Q(s,a) \leftarrow \Sigma_{s', r'}p(s', r'|s, a)[r + \gamma \Sigma_{a'}\pi(a'|s')Q(s', a')]$\\
\indent \indent \indent  $\delta \leftarrow max(\Delta|\textit{q} - Q(s,a)|)$\\
\indent until $\Delta < \theta$(a small positive number determining the accuracy of estimation) \\

3. Policy Improvement\\

$\textit{policy-stable} \leftarrow \textit{true}$\\
\indent For each $s \in S:$\\
\indent \indent $\textit{old-action} \leftarrow \pi_{s} \\$
\indent \indent $\pi(s) \leftarrow argmax_{a}\Sigma_{s', r'}p(s', r'|s, \pi(s,a))[r + \gamma \Sigma_{a'}\pi(a'|s')Q(s', a')]$\\
\indent \indent If old-action $\neq \pi(s,a)$, the policy-stable $\leftarrow \textit{false}$\\
\indent If $\textit{policy-stable}$, then stop and return $Q \approx Q_{*}$ and $\pi \approx \pi_{*}$; else go to 2.
\\\\
\textit{Exercise 4.6} Suppose you are restricted to considering only policies that are $\epsilon$-soft,
meaning that the probability of selecting each action in each state, s, is at least $\frac{\epsilon}{|A(s)|}$.
Describe qualitatively the changes that would be required in each of the steps 3, 2, and 1,
in that order, of the policy iteration algorithm for $v_{*}$ on page 80. 
\\

Every action now has a minimum chance of being selected given by $\frac{\epsilon}{|A(s)|}$. 

Policy is now stochastic. The policy can be arbitrarily initialized as

before, but control statements must be implemented so as to guarantee the 

minimum probability for all actions and account for this on other actions so

that $P_{total} = 1$ still. The new policy evaluation and improvement will be:\\

2. Policy Evaluation\\

\indent Loop:\\
\indent \indent $\Delta \leftarrow 0$\\
\indent \indent Loop for each s $\in$ S:\\
\indent \indent \indent $\textit{v} \leftarrow V(s)$\\
\indent \indent \indent $V(s) \leftarrow \Sigma_{a}\pi(a|s)\Sigma_{s', r'}p(s', r'|s, a)[r + \gamma V(s')]]$\\
\indent \indent \indent  $\delta \leftarrow max(\Delta|\textit{V} - V(s)|)$\\
\indent until $\Delta < \theta$(a small positive number determining the accuracy of estimation) \\

3. Policy Iteration\\

$\textit{policy-stable} \leftarrow \textit{true}$\\
\indent For each $s \in S:$\\
\indent \indent $\textit{old-action} \leftarrow \pi_{s}$ \\
\[ \pi(s) = \begin{cases} 
          $1 - \epsilon + \frac{\epsilon}{|A(s)|}$, if $a = argmax_{a}$V(s)\\$
          $\frac{\epsilon}{|A(s)|}$, otherwise
       \end{cases}
    \]
\indent \indent If old-action $\neq \pi(s)$, the policy-stable $\leftarrow \textit{false}$\\
\indent If $\textit{policy-stable}$, then stop and return $V \approx V_{*}$ and $\pi \approx \pi_{*}$; else go to 2.\\

In words, when we find the argmax(a) for state s, we no longer can select 

a with probability 1. Instead, we must assign it the max probability under 

our $\epsilon$-soft constraint. For example, if there are five states total, $\frac{\epsilon}{5}$ must be 

the minimum probability assigned to each state. Now once argmax(a) is 

found and we know p(a) is no longer 1, and since we still want to maximize 

p(a), we subtract: 1 - $\epsilon$ + $\frac{\epsilon}{5} = 1 - \frac{4}{5}\epsilon$. $\frac{4}{5}\epsilon$ is the total probability of actions 

not including argmax(a) under our $\epsilon$-soft constraint.
\\\\
\textit{Exercise 4.8} Why does the optimal policy for the gambler’s problem have such a curious
form? In particular, for capital of 50 it bets it all on one flip, but for capital of 51 it does
not. Why is this a good policy?. 
\\

$P_{h} = 0.4$, so if you bet $1.00$ each round, you are expected to slowly lose. 

Hence, the $50.00$ mark acts as a threshold: if you are at $50.00$, you bet 

everything to minimize the amount of rounds played to win and thus 

maximize the probability of winning. Now if you are above $50.00$ but below 

$75.00$, you have the threshold of $50.00$ to fall back on. You thus want to

stake anything that you could reach to $75.00$ with but keep you at or above 

$50.00$, which is the reason for the small triangular "tics", which is the 

class of optimal policies referred to in the above paragraph (for example, 

at $65$, you could arbitrarily stake $10$ or $15$). For the spike at $75$, 

the logic is similar; we can win by betting $25.00$, but if we lose we fall 

back to $50$ (and if we are above $75$ we can stake arbitrarily to win, but 

if we lose keep ourselves at or above $75$). Likewise for $25$: we bet $25.00$ 

on the chance we jump to $50.00$ and then bet $50.00$ if we make it (and again, 

the tics are because if we are under $25$ we can arbitrarily stake to get up to 

$25$, and if we are above $25$ but below $50$, we can arbitrarily stake to get 

up to $50$ but stay at or above $25$ if we lose). To drive the point home again, 

we want to minimize the expected number of rounds played to win.
\\\\
\textit{Exercise 4.10} What is the analog of the value iteration update (4.10) for action values,
q k+1 (s, a)?
\\

$q_{k+1}(s, a) = \Sigma_{s', r'}p(s', r'|s, a)[r + \gamma max_{a'}q_{k}(s', a')]$

\maketitle

\section{Introduction}

\end{document}
